\section{Full common-space extension with finitely many retained boundary zeros}
\label{sec:finite-boundary-full-matrix}

The necessary boundary condition in MA-MO-025/026 has a matching
sufficient case. A finite boundary divisor can be treated by its
actual right boundary distributions, provided the remaining zeros
have a uniform interior margin. The polynomial weight then absorbs
the boundary poles at exactly the previously established jet
threshold. The construction extends the original wide-inverse
source matrix; it does not replace that inverse by an arbitrary
principal value.

Use either of the following already fixed source settings:
the nontrivial finite-order character sector of MA-MO-022, with
its number field, coherent Fourier normalizations and structural
place set (MA-MO.470); or the rational principal sector
\(K=\mathbb Q,S=\{\infty\}\) and the actual polar completion
of MA-MO-016/018. In the first case let
\(\omega\in\{\eta,\bar\eta\}\); in the second case there is
one self-dual multiplier. Write
\[
 \begin{gathered}
 0<a<1/2,\qquad s\ge0,\qquad r=s/2,\\
 H_{\pm a,s}
   =L^2\bigl(e^{\pm2ax}(1+x^2)^{s/2}\,dx\bigr),\\
 H=H_{\Sigma,a,s}=H_{a,s}\cap H_{-a,s},\qquad
 \|y\|_H^2=\|y\|_{H_{a,s}}^2+\|y\|_{H_{-a,s}}^2 .
 \end{gathered}
 \tag{MA-MO.700}
\]
The function \(\phi\) is the original real smooth step,
equal to one far left and zero far right, with
\(\phi'\in C_c^\infty\). Put \(q=1-\phi\). No symmetry
or Gevrey condition is imposed.

The new arithmetic hypothesis is the following finite-boundary
condition on each completed centered divisor:
\[
 \begin{gathered}
 |\Re\lambda|\le a\quad\hbox{for every completed zero }\lambda,\\
 \{\lambda:|\Re\lambda|=a\}\ \hbox{is finite},\\
 \text{for some }\delta>0,\qquad
 |\Re\lambda|<a\ \Longrightarrow\
                       |\Re\lambda|\le a-\delta .
 \end{gathered}
 \tag{MA-MO.701}
\]
Here \(\lambda=\rho-1/2\), with its full multiplicity
\(m_\omega(\lambda)\). The condition is required for both
dual characters in the nontrivial case. Shrinking \(\delta\)
allows \(0<\delta<a\). No claim that (MA-MO.701) holds
arithmetically is made.

\paragraph{MA-MO-027 (finite-boundary sufficiency in the actual norm).}
Assume (MA-MO.701) and
\[
        s>2m_\omega(\lambda)-1
               \quad\hbox{at every boundary zero }|\Re\lambda|=a .
 \tag{MA-MO.702}
\]
The condition is vacuous if the boundary divisor is empty.
For every \(g\in C_c^\infty(\mathbb R)\), the original
smoothed matrix \(C_g(P_+-P_-)\) has a unique trace-class
extension on the actual completion:
\[
 \begin{array}{ll}
 \mathscr D_{\eta,g}:H\longrightarrow H,
       &\text{nontrivial character sector},\\[2pt]
 \mathscr D_g:\mathcal W_{a,s}\longrightarrow\mathcal W_{a,s},
       \quad\mathcal W_{a,s}=H\oplus\mathbb C^2,
       &\text{rational principal sector}.
 \end{array}
 \tag{MA-MO.703}
\]
On each fixed compact test support the extensions depend
continuously on \(g\) in trace norm. The two polar coordinates
in the second line are the actual source-section coordinates,
not extra bare exponentials inserted into \(H\).

In the nontrivial case \(J_\eta=\overline{I_\eta}^{\,H}\)
is invariant, and the induced quotient action is \(-C_g\).
Writing \(B_{\eta,g}^{\partial}\) for the actual restriction
to this closed ideal, one has
\[
 \begin{aligned}
 \operatorname{Tr}_H\mathscr D_{\eta,g}
   &=\operatorname{Tr}_{J_\eta}B_{\eta,g}^{\partial}
                  -\sum_\rho m_\eta(\rho)G(\rho-1/2),\\
 \det_H(1-z\mathscr D_{\eta,g})
   &=\det_{J_\eta}(1-zB_{\eta,g}^{\partial})
                         \det_{H/J_\eta}(1+zC_g),
               \qquad G=\mathcal Mg .
 \end{aligned}
 \tag{MA-MO.704}
\]
This statement supplies the actual boundary extension and
its invariant-subspace identities. It does not identify
the first trace in (MA-MO.704) with a local scalar formula
by silently reusing a no-boundary inverse estimate.

Within the class (MA-MO.701), condition (MA-MO.702) is
also necessary for any nonzero compact smooth test to have
a bounded original full-matrix extension: this is exactly
MA-MO-025 for the nontrivial sector and MA-MO-026 for the
rational polar completion. Thus it characterizes bounded
full-matrix existence in this finite-boundary, uniform-margin
case. It is not a sufficiency theorem for an infinite divisor
approaching an edge without a uniform margin.

\begin{proof}
\emph{Boundary reciprocals and the exact weight threshold.}
Use the original forward multipliers
\[
 A_\omega(z)=
 \begin{cases}
 E_{\omega,S}(1/2+z)L_K(1/2+z,\omega),
                           &\text{nontrivial case},\\
 \zeta(1/2+z),              &\text{rational case},
 \end{cases}
 \qquad j_{\omega,\epsilon}(u)=A_\omega(a+\epsilon+iu)^{-1}.
 \tag{MA-MO.705}
\]
The finite Euler factors have their zeros on \(\Re z=-1/2\).
Choose \(0<\epsilon_0<1/2-a\) sufficiently small. A narrow
closed strip about the line \(\Re z=a\) contains no zeros
except the finitely many right boundary zeros
\(\lambda_j=a+i\gamma_j\). It also avoids the rational
pole at \(z=1/2\).

In a disk around \(\lambda_j\), write the genuine Laurent
principal part
\[
 A_\omega(z)^{-1}
       =\sum_{k=1}^{m_j}c_{j,k}(z-\lambda_j)^{-k}
                               +\text{holomorphic function},
               \qquad c_{j,m_j}\ne0 .
 \tag{MA-MO.706}
\]
Choose compact smooth functions \(\chi_j(u)\), equal to
one near \(\gamma_j\), with sufficiently small disjoint
supports. Subtracting these localized principal parts gives
\[
 j_{\omega,\epsilon}(u)=r_{\omega,\epsilon}(u)
    +\sum_j\sum_{k=1}^{m_j}
       c_{j,k}\chi_j(u)(\epsilon+i(u-\gamma_j))^{-k},
                   \qquad 0<\epsilon\le\epsilon_0 .
 \tag{MA-MO.707}
\]
Every derivative of \(r_{\omega,\epsilon}\) has a polynomial
bound in \(u\), uniform for these \(\epsilon\), and these
functions converge to a smooth \(r_{\omega,0}\) with the
same weighted derivative bounds. More precisely, for each
derivative there is a fixed inverse polynomial weight in
which this convergence is uniform; a further derivative
in the real part and the mean value theorem give this
assertion. Here is the exact
minimum-modulus input. Factoring the finitely many zeros
in (MA-MO.706) leaves a multiplier without zeros in a
fixed neighborhood of the positive line, with a fixed
horizontal margin to every remaining zero. The
fixed-gap proof of MA-RS-004, already used for each fixed
Hecke character in MA-MO-012/022, gives polynomial
reciprocal derivative bounds outside the finitely many
disks. The factors removed are polynomials, so preserve
such bounds. Inside the disks the Laurent remainders
are holomorphic on a slightly larger compact disk;
Cauchy estimates are uniform. The transition regions
of the \(\chi_j\) stay a positive distance from the
poles. These three regions prove the asserted estimates
and convergence. In the rational case the distant pole
of \(A\) is a zero of its reciprocal; it creates no
additional inverse singularity.

For each positive integer \(k\), the elementary Laplace
identity is
\[
 (\epsilon+iv)^{-k}
     =\frac1{(k-1)!}\int_0^\infty
                 t^{k-1}e^{-\epsilon t}e^{-ivt}\,dt .
 \tag{MA-MO.708}
\]
With the unitary plus Fourier transform, its transform
is \(\sqrt{2\pi}\,t^{k-1}e^{-\epsilon t}
\mathbf1_{t>0}/(k-1)!\). Consequently it converges as
\(\epsilon\downarrow0\) in \(H^{-r}(\mathbb R)\) if
and only if
\[
       \int_0^\infty t^{2k-2}(1+t^2)^{-r}\,dt<\infty,
               \qquad\text{equivalently }r>k-\tfrac12 .
 \tag{MA-MO.709}
\]
Necessity follows from the same nonnegative integral,
and sufficiency and convergence follow by dominated
convergence. Translation and multiplication by a fixed
compact smooth function are bounded on \(H^{-r}\).
Thus (MA-MO.702) gives the actual boundary distribution
\[
 j_{\omega,+}
    =r_{\omega,0}+
       \sum_{j,k}c_{j,k}\chi_j
             (0+i(\,\cdot-\gamma_j))^{-k}.
 \tag{MA-MO.710}
\]
Each singular summand has compact support and belongs to
\(H^{-r}\). Its sign and delta-derivative contributions
are fixed by the limit from \(\Re z>a\). No alternative
principal-value prescription has been chosen.

\emph{Identification with the original inverse before using a kernel.}
For \(y\in\mathcal E\), let \(F_{\omega,y}\) be the
original wide inverse series (MA-MO.477), or (MA-MO.323)
in the rational case. That series is locally absolutely
convergent with all derivatives. On every line \(b>a\)
its Mellin transform is the original wide quotient shifted
without crossing any zero:
\[
 \begin{gathered}
 e^{bx}F_{\omega,y}\in\mathscr S(\mathbb R),\qquad
 V_bF_{\omega,y}=j_{\omega,b-a}V_by,\qquad
 V_b=\mathcal F_+M_{e^{bx}},\\
                         a<b<a+\epsilon_0 .
 \end{gathered}
 \tag{MA-MO.711}
\]
Indeed such a line has a positive gap from every zero;
the already proved wide-inverse contour argument of
MA-MO-018/022 applies, with \(a_*=a<b\).
The same argument supplies every fixed right weight
\(b>a\), including arbitrarily large ones.

The transforms \(V_by\) converge in Schwartz space to
\(V_ay\). Equations (MA-MO.707)--(MA-MO.710) therefore
give convergence in tempered distributions, and
Fourier inversion gives
\[
        V_aF_{\omega,y}=j_{\omega,+}V_ay
                    \quad\hbox{as tempered distributions}.
 \tag{MA-MO.712}
\]
This is an identity with the original function: on every
compact \(x\)-interval,
\(e^{bx}F_{\omega,y}\to e^{ax}F_{\omega,y}\) smoothly,
so the distributional limit just obtained has that
local value. The compact singular distributions in
(MA-MO.710), after multiplication by \(V_ay\), have
finite order. Their inverse Fourier transforms, and
all derivatives, have at most polynomial growth.
The regular part times \(V_ay\) is Schwartz. This
also proves the needed polynomial growth of
\(e^{ax}F_{\omega,y}\), rather than assuming it
from a formal boundary integral.

Multiplying (MA-MO.707) by \(A_\omega(a+\epsilon+iu)\)
and taking the same distributional limit proves
\[
                    A_\omega(a+iu)j_{\omega,+}(u)=1 .
 \tag{MA-MO.713}
\]
The multiplication is legitimate because the forward
factor and all its derivatives have polynomial bounds;
on each compact singular support the factors converge
in every smooth norm. This identity is established
before any diagonal commutator cancellation.

Finally \(qF_{\omega,y}\in\mathcal E_\omega\):
it vanishes far left and has all right weighted
Schwartz estimates from (MA-MO.711).
For a nontrivial sector, the structural-\(S\)
inclusion (MA-MO.476) puts it in the actual local
\(\mathcal T_\omega\), with its source topology.
For the rational sector,
\[
       u_y(t)=\frac{q(\log|t|)F_y(\log|t|)}
                         {2\sqrt{|t|}},\qquad u_y(0)=0
 \tag{MA-MO.714}
\]
is the actual even additive Schwartz function,
zero near additive zero. Thus both subsequent
source Fourier transforms are on their original domains.

\emph{A trace-class lemma for the singular input distributions.}
We give the norm argument needed at the boundary.
If \(d\in H^{-r}(\mathbb R)\) has compact support
and \(L\in\mathscr S(\mathbb R^2)\), then
\[
       T_{d,L}f(t)=\langle d(u),L(t,u)f(u)\rangle
          \quad\hbox{defines }\quad
       T_{d,L}\in\mathcal S_1(H^r,H^R)
                    \quad\hbox{for every fixed }R\ge0 .
 \tag{MA-MO.715}
\]
The pairing is bilinear distributional pairing;
it is not a pointwise multiplication of two
undefined distributions. To prove the assertion,
expand in the normalized Hermite functions of
the output variable:
\(L(t,u)=\sum_{n\ge0}h_n(t)L_n(u)\).
Repeated use of the harmonic oscillator in \(t\)
shows that every Schwartz seminorm of \(L_n\)
decreases faster than every power of \(n+1\).
For a smooth multiplier \(v\), the Fourier
convolution formula and
\((1+|\xi+\eta|^2)^{r/2}
 \le C_r(1+|\xi|^2)^{r/2}(1+|\eta|^2)^{r/2}\)
give
\[
 \|vf\|_{H^r}
   \le C_r\bigl\|(1+|\xi|^2)^{r/2}\widehat v(\xi)
                      \bigr\|_{L^1}\|f\|_{H^r}.
 \tag{MA-MO.716}
\]
Hence the norms of
\(\ell_n(f)=\langle d,L_nf\rangle\) decrease rapidly.
The \(H^R\)-norms of \(h_n\) grow at most polynomially:
integer derivative norms follow from the two Hermite
ladder identities, and a larger integer derivative
norm bounds the fractional one. It follows that
\(\sum_n\|h_n\|_{H^R}\|\ell_n\|<\infty\).
This is an explicit nuclear series between Hilbert
spaces and proves (MA-MO.715).

The proof also proves continuity in \(d\) with its
\(H^{-r}\) norm and in finitely many Schwartz
seminorms of \(L\). In particular \(d_\epsilon\to d\)
in \(H^{-r}\) and \(L_\epsilon\to L\) in Schwartz
space imply convergence in trace norm. A smooth
polynomially bounded input factor times \(L\) is
an ordinary Schwartz kernel and has the same
conclusion. Finally \(V_a\) identifies the actual
weighted norm with \(H^r\):
\[
           \|y\|_{H_{a,s}}
                =\|(1-\partial_t^2)^{r/2}V_ay\|_2,
                          \qquad r=s/2 .
 \tag{MA-MO.717}
\]
Thus the lemma concerns exactly the specified
polynomial weight.

\emph{Cancellation before the boundary pairing.}
Put \(P(\xi)=\int e^{\xi x}\phi'(x)\,dx\).
For \(b=a+\epsilon\), let \(z=b+it,w=b+iu\).
The ordinary commutator kernel (MA-MO.184) can
be written, with its sign unchanged, as
\[
 \begin{aligned}
 k_{\omega,b,g}(t,u)&=L_{\omega,b,g}(t,u)
                                      j_{\omega,\epsilon}(u),\\
 L_{\omega,b,g}(t,u)
   &=\frac{G(z)P(z-w)}{2\pi}
                          \frac{A_\omega(z)-A_\omega(w)}{z-w}.
 \end{aligned}
 \tag{MA-MO.718}
\]
The divided difference has its removable ordinary
value at \(z=w\), before the singular reciprocal
is paired with anything. The equality follows from
\(A(z)[j(z)-j(w)]=(A(w)-A(z))j(w)\) and
\(z-w=i(t-u)\). Equivalently it uses (MA-MO.713)
in the limit. There is no product of a diagonal
principal-value kernel with an independent pole
distribution.

The functions \(L_{\omega,b,g}\) converge in
\(\mathscr S(\mathbb R^2)\) as \(b\downarrow a\).
To check all derivatives, the forward multiplier
has polynomial derivative bounds on this closed
strip. Its divided difference has polynomial
derivative bounds too, by integrating its derivative
along the vertical segment from \(w\) to \(z\).
The factors \(G(b+it)\) and \(P(i(t-u))\) supply
rapid decrease in the independent directions
\(t\) and \(t-u\). Equations (MA-MO.707),
(MA-MO.710) and (MA-MO.715) therefore give a
trace-class limit \(K_{\omega,a}^{\partial}(g)\)
on \(H_{a,s}\).

The crossed Fourier kernel has the same type of
controlled singularity, now with the \emph{dual}
reciprocal:
\[
 \begin{aligned}
 o_{\eta,b,g}(t,u)
   &=L^{\mathrm{cross}}_{\eta,b,g}(t,u)
                                      j_{\bar\eta,\epsilon}(u),\\
 L^{\mathrm{cross}}_{\eta,b,g}(t,u)
   &=\frac{G(z)A_\eta(z)m_\eta(z)}{2\pi}
                              \frac{P(-z-w)}{-z-w}.
 \end{aligned}
 \tag{MA-MO.719}
\]
In the rational case omit the subscripts on the
self-dual factors. This is exactly (MA-MO.481)
or (MA-MO.331) on the line \(b>a\), not a
new Fourier normalization. The real part of
\(-z-w\) stays at most \(-2a<0\).
The actual local multiplier \(m_\eta\), including
the finite phases, conductor and different, is
holomorphic with polynomial derivative bounds
on the compact positive strip under consideration.
The rapid directions are \(t\) and \(t+u\).
Thus \(L^{\mathrm{cross}}_{\eta,b,g}\) also
converges in Schwartz space, and (MA-MO.715)
gives an actual trace-class limit
\(O_{\eta,a}^{\partial}(g):H_{a,s}^{\bar\eta}
\to H_{a,s}^{\eta}\).

\emph{Actual compatible character rows.}
The third term needed in the first row is
\(N_a(g)=C_gM_\phi:H\to H_{a,s}\).
It is trace class by the exact physical kernel
(MA-MO.332): \(x-y\) has compact support,
the normalized weight ratio decreases exponentially
as \(y\to-\infty\), and \(\phi(y)=0\) far right.
Every derivative obeys the same estimates.
This argument uses the sum norm on the input;
it makes no trace-class assertion for an unchanged
one-sided half-line multiplication.

Define
\[
 \mathscr R_{\eta,a}^{\partial}(g)y
       =K_{\eta,a}^{\partial}(g)y
                     -N_a(g)y+O_{\eta,a}^{\partial}(g)Ry .
 \tag{MA-MO.720}
\]
It is trace class from \(H\) to \(H_{a,s}\).
For \(y\in\mathcal E\), it is the first row of
the original source matrix. To prove this exact
agreement, first use the already proved source
identity (MA-MO.479)--(MA-MO.482) at every
\(b>a\). Both inverse functions there are the
fixed original functions identified in
(MA-MO.711). Their cutoff vectors are actual
local source vectors by (MA-MO.714) and the
structural inclusion, and their theta outputs
are in \(\mathcal E\). In the Fourier picture
the input \(V_by\) converges in \(H^r\), the
two kernel operators converge in trace norm by
the preceding construction, and the transform
of that fixed actual output converges to its
\(a\)-line transform. The half-line term has
the same direct limit. Thus the identity passes
to (MA-MO.720). This argument takes limits
on the genuine smooth core; it does not assert
that identity maps between all varying-width
Hilbert spaces are bounded.

The second row is
\(R\mathscr R_{\bar\eta,a}^{\partial}(g^\vee)R\).
On \(\mathcal E\) the two functions agree by
the actual dual Fourier relation and nonpolar
Poisson summation (MA-MO.479). The paired map
is trace class into \(H_{a,s}\oplus H_{-a,s}\).
The diagonal image of \(H\) is closed and
\(\mathcal E\) is dense in \(H\), so the range
of the paired map is in that diagonal on the
whole completion. Its inverse diagonal isometry
gives the first operator in (MA-MO.703),
with agreement on the entire original smooth
nontrivial common sector.

\emph{The rational cutoff moments and polar rows.}
The rational case requires two further checks.
For \(y\in\mathcal E\), the original \(F_y\)
has polynomially weighted \(a\)-line growth
at the left by (MA-MO.712), and arbitrary
right weighted decrease. Therefore its Mellin
integral at \(1/2\) is absolutely convergent.
The wide inverse quotient continues through
that point and has a zero there, so
\[
          \mathcal MF_y(1/2)=0 .
 \tag{MA-MO.721}
\]
This is the original zero total moment, not
a moment imposed on the input.

Set \(p_a(x)=\phi(x)e^{(1/2-a)x}\).
It is Schwartz, since \(a<1/2\) and
\(\phi\) is zero far right. With the unitary
plus Fourier normalization define
\[
 \begin{aligned}
 \lambda_0(y)
    &=\big\langle j_+(u),
                     \widehat p_a(-u)(V_ay)(u)\big\rangle,\\
 \lambda_1(y)&=\lambda_0(Ry).
 \end{aligned}
 \tag{MA-MO.722}
\]
The regular summand in (MA-MO.710) times
\(\widehat p_a(-u)\) is Schwartz.
Each compact singular summand is in \(H^{-r}\),
and multiplication by \(\widehat p_a(-u)\)
is bounded on \(H^r\), by (MA-MO.716).
Thus \(\lambda_0\) is bounded on \(H_{a,s}\),
\(\lambda_1\) is bounded on \(H_{-a,s}\), and
both are bounded on the actual sum space.

For \(y\in\mathcal E\), bilinear Parseval
on the line \(b>a\), followed by the same
boundary limit, gives exactly
\[
 \lambda_0(y)=\int\phi(x)e^{x/2}F_y(x)\,dx,
              \qquad \lambda_1(y)=\lambda_0(Ry).
 \tag{MA-MO.723}
\]
There is no extra Fourier constant: both
transforms in (MA-MO.722) are unitary.
The physical integral converges by the
left estimate proved before (MA-MO.721);
on lines \(b>a\) its weight is
\(\phi(x)e^{(1/2-b)x}\), whose Fourier
transform converges in Schwartz space to
\(\widehat p_a\).

The actual source \(h_y=u_y+\mathfrak F u_{Ry}\)
from (MA-MO.326) consequently has
\[
       h_y(0)=-\lambda_1(y),\qquad
       \int h_y=-\lambda_0(y).
 \tag{MA-MO.724}
\]
Indeed \(u_y(0)=0\), while its integral is
\(\int e^{x/2}qF_y=-\lambda_0(y)\), by
(MA-MO.721); Fourier interchanges evaluation
and integration. These statements use the
actual additive Schwartz functions (MA-MO.714).

Retain the fixed Gaussian sections \(f_0,f_1\)
and vectors \(e_0=\Theta f_0\),
\(e_1=\Theta f_1-e^{-x/2}\) of
(MA-MO.333), with \(Re_0=e_1\).
They belong to \(H\) since \(a<1/2\).
After subtracting the actual polar sections,
the first row and the polar coordinates are
\[
 \begin{aligned}
 T_{a,g}^{\partial}y
   ={}&K_a^{\partial}(g)y-N_a(g)y
                                 +O_a^{\partial}(g)Ry\\
      &+G(-1/2)\lambda_1(y)e_0
                     +G(1/2)\lambda_0(y)e_1,\\
 \nu_g(y)
   ={}&\bigl(-G(-1/2)\lambda_1(y),
                         -G(1/2)\lambda_0(y)\bigr).
 \end{aligned}
 \tag{MA-MO.725}
\]
The signs are those of the actual centered residue
\(\Theta h=Z_bF_h+(\int h)e^{-x/2}\)
in (MA-MO.334). That direct theta contour
shift crosses only the pole of zeta, not a pole
at any of its zeros. Alternatively, all
identities (MA-MO.328)--(MA-MO.335) hold at
\(b>a\); the now proved moment and kernel
limits give exactly (MA-MO.725).
Every term in its first line is trace class
\(H\to H_{a,s}\), and \(\nu_g\) is bounded
finite rank.

The reflected second row is
\(RT_{a,g^\vee}^{\partial}R\).
It agrees on \(\mathcal E\) by
\(h_{Ry}=\mathfrak Fh_y\) and the interchange
of the two Gaussian sections. As in the
character case, the closed diagonal gives
a unique trace-class \(T_g^\partial:H\to H\).
The negative source column is therefore
\(y\mapsto(T_g^\partial y,\nu_g(y))\).
The two remaining columns \(C_gDs_j\) are
fixed actual smooth common-sum vectors,
already constructed without a zero-location
assumption in MA-MO-016. Appending those
finite columns gives
\[
 \mathscr D_g(y,c)
       =(T_g^\partial y,\nu_g(y))
                           +\sum_{j=0}^1c_jC_gDs_j .
 \tag{MA-MO.726}
\]
This is trace class on \(\mathcal W_{a,s}\)
and agrees with the original matrix on the
dense source \(\mathfrak s(\mathcal E\oplus\mathbb C^2)\).
This proves the rational part of (MA-MO.703).

\emph{The quotient statements and the scope of the result.}
In the nontrivial case the exact smooth formula
has \(P_+\mathcal E\subset I_\eta\).
It also preserves \(I_\eta\) after subtracting
the identity and convolving. The bounded
extension therefore preserves \(J_\eta\), and
\(\mathscr D_{\eta,g}+C_g\) has range in
\(J_\eta\) by density. Hence its quotient is
\(-C_g\). Condition (MA-MO.702) retains all
boundary jets, by the exact inequality
\(2j<s-1\) of MA-WP-001--003: for a
multiplicity \(m\), every \(0\le j<m\)
is retained precisely when \(s>2m-1\).
All interior zeros are retained as well.
The actual quotient trace is consequently
the full absolutely convergent sum in
(MA-MO.704). Triangular compression on
\(J_\eta\oplus J_\eta^\perp\), followed by
finite-rank trace-norm approximation, proves
both the trace and Fredholm determinant
identities there.

The source core is dense in each claimed
completion, proving uniqueness. Each
Schwartz kernel estimate and each finite
distribution pairing uses only finitely
many derivatives of \(g\) on a fixed
compact support. The finite source columns
have the same continuity. Thus the maps
are continuous into trace class.
Finally MA-MO-025/026 already refutes
boundedness for every nonzero compact
test when a boundary multiplicity violates
(MA-MO.702). This proves the stated exact
characterization within (MA-MO.701).
\end{proof}

\paragraph{Remaining scope and framework placement.}
The proof uses the finite edge divisor twice:
to have only finitely many \(H^{-s/2}\) input
singularities, and to retain uniform polynomial
estimates for the remaining reciprocal.
It gives no summability estimate for infinitely
many approaching interior poles, and does not
replace their actual lifting costs by their
formal multiplicities. The broad full-retention
sufficiency question of MA-MO-025 remains
separate.

The rational operator is an endomorphism of
the polar completion \(\mathcal W\), not
automatically of \(\mathcal W_+\); the
rank-two outgoing obstruction of MA-MO-023
is unchanged. In the nontrivial sector the
ideal restriction in (MA-MO.704) is actual
and trace class, but a local evaluation of
that restriction is not assumed merely
from the boundary construction.

The primary source maps and structural
interfaces are the already checked Meyer
locations in MA-MO-018/022; the present
proof introduces the boundary-distribution
estimate, not a new reading of those
sources. In the inherited cote~1 framework,
scans~079--080 and~111--112 through
ALPOGE\_BRIDGES BR-TR-01/02, a boundary
multiplicity is a finite-length datum,
whereas (MA-MO.709), (MA-MO.715) and
the actual source comparison supply the
additional analytic continuity and
trace-ideal structure. No arithmetic
positivity or new zero-location theorem
is inferred from the finite-length class.
